\documentclass[11pt,a4paper]{article}
\usepackage[utf8]{inputenc}
\usepackage[T1]{fontenc}
\usepackage[spanish]{babel}
\usepackage{hyperref}
\usepackage{geometry}
\geometry{margin=2.5cm}
\usepackage{parskip}

\title{329 sabe lo que cuesta, y puede corregirlo}
\author{Alejandro Valenzuela Rivera\\\small Universidad Autónoma de Nuevo León}
\date{\today}
\begin{document}
\maketitle

\begin{quote}
\textbf{Sesgo declarado.} Este trabajo se sostiene desde el sesgo declarado de su autor, y no finge neutralidad. Lo que afirma se puede tumbar, y ese es el punto: el falsificador está escrito, no escondido.
\end{quote}

\section*{Resumen}
Entre Terminator y WALL-E falta el punto medio: el que sabe lo que cuesta lo que hizo y aún puede corregirlo.

\section{Por qué se sostiene}
La regla se sostiene mientras alguien la pueda contestar con lo que Terminator, WALL-E y el punto medio muestra.

\section{De dónde viene}
Tesis/publicable/casos.md.

\section{El caso}
Terminator: no sabe lo que cuesta. WALL-E: sabe y no puede cambiarlo.

\section{Método}
Sin registro del efecto no hay corrección: el que no sabe lo que costó no puede deshacerlo.

\section{Límites, y qué lo tumba}
Un caso donde saber el costo haga imposible corregir.

\section{Lo que queda abierto}
\textquestiondown{}Sabe lo que costó, y puede corregirlo?

\section*{Cómo se leyó esta tesis}
Esto se escribió en el repositorio man-made-empathy, del que este candado es una pieza: candados/06 del caso/329 sabe lo que cuesta, y puede corregirlo.md. La máquina que lo compuso lo hizo desde este mismo archivo, y no agregó nada que no estuviera escrito. Si algo de esto es falso, se dice en el mismo sitio, y se corrige con fecha.

\end{document}
